\section{Conclusion}
\label{sec:conclusion}

We tested five claims about attention spectra over training, with per-head tracking across 68 runs and 120 information runs. Heads do not split into high- and low-rank populations. Under standard initialization they start near-uniform and gain rank, so rank collapse is not a universal training dynamic. The bound $I(A;X)\le H(Y|X)+O(1/\sqrt n)$ is false, and our variance bound is a valid replacement. The $0.3\log n$ pruning rule is unsafe, though entropy carries pruning information on \cifar. One finding survives in modified form: the rank fall from peaked starts grows strongly with depth.

{\bf The takeaway:} the initial state, not training alone, decides whether attention rank falls.

Future work should sweep query and key gain and softmax temperature to map how the start state carries through depth, test a per-layer entropy ranking on larger vision models, ask whether stable spatial roles explain the link between entropy and importance on vision, and analyze per-input distributions, since per-head means may hide input-level bimodality. If a derivation of the H4 bound exists, its assumptions deserve a direct test.
